\documentclass[11pt,a4paper]{article}

\usepackage[utf8]{inputenc}
\usepackage[T1]{fontenc}
\usepackage[brazil]{babel}
\usepackage[a4paper,margin=2.4cm]{geometry}
\usepackage{listings}
\usepackage{xcolor}
\usepackage{enumitem}
\usepackage[hidelinks]{hyperref}
\usepackage{xurl}
\usepackage{seqsplit}
\usepackage{etoolbox}

\hypersetup{
  colorlinks=true,
  linkcolor=blue!45!black,
  urlcolor=blue!55!black,
  citecolor=black,
}

% Blocos de codigo NAO podem ser cortados entre paginas: cada lstlisting vai
% dentro de um minipage indivisivel, entao o bloco inteiro desce para a proxima
% pagina se nao couber, em vez de partir no meio.
\BeforeBeginEnvironment{lstlisting}{\par\noindent\begin{minipage}{\linewidth}}
\AfterEndEnvironment{lstlisting}{\end{minipage}\par}

% Nao esticar o espaco vertical para encher a pagina: quando um bloco desce
% inteiro, o vazio fica no fim da pagina, sem abrir buracos entre as secoes.
\raggedbottom

\definecolor{codebg}{gray}{0.96}
\lstset{
  basicstyle=\ttfamily\footnotesize,
  backgroundcolor=\color{codebg},
  breaklines=true,
  columns=fullflexible,
  keepspaces=true,
  frame=single,
  rulecolor=\color{gray!40},
  framesep=5pt,
  xleftmargin=6pt,
  xrightmargin=6pt,
  showstringspaces=false,
  aboveskip=8pt,
  belowskip=8pt,
}

\setlist{itemsep=2pt,topsep=3pt}

\title{Infraestrutura de Python da AURORA}
\author{NIPS-CERN}
\date{\today}

\begin{document}

\maketitle

\begin{abstract}
\noindent
Este documento descreve, de forma geral, como o Python funciona dentro da
AURORA: qual interpretador embarcamos, de onde ele vem (o bundle MSYS2), quais
bibliotecas estão de fato instaladas, como um testbench cocotb roda de ponta a
ponta, onde fica cada componente com os nomes reais, e o plano para distribuir
bibliotecas sob demanda, com o Python do próprio usuário via terminal TCMD como
saída para as bibliotecas que não empacotarmos.
\end{abstract}

% Corpo alinhado a esquerda: nomes de arquivo e caminhos longos em \texttt nao
% escapam para fora da margem porque as linhas nao sao esticadas ate a direita.
\raggedright

\section{Visão geral}

O Python na AURORA existe para uma finalidade específica: rodar testbenches
cocotb. cocotb é um framework em Python para verificação de hardware; em vez de
escrever o testbench em Verilog, o usuário escreve um arquivo \texttt{.py} que
estimula e checa o DUT enquanto um simulador (Icarus Verilog ou Verilator) roda
o design por baixo.

Por isso o nosso Python é um runtime enxuto, embarcado, com um único trabalho.
Ele não é um Python de propósito geral com muitas bibliotecas; é um interpretador
mínimo que carrega o cocotb. Ele vive dentro do mesmo bundle MSYS2/mingw64 que
traz o resto da toolchain (iverilog, vvp, verilator, yosys, perl, g++, make), de
modo que há um único Python para os dois fluxos cocotb (Icarus e Verilator).

Os pontos gerais:

\begin{itemize}
  \item Interpretador embarcado: \texttt{python.exe} dentro do bundle
        \texttt{components/Packages/msys/mingw64/bin/}. É um build MinGW do MSYS2,
        não o CPython oficial (MSVC) da python.org. Essa distinção importa para o
        plano de bibliotecas (Seção 8).
  \item Versões fixadas: Python 3.12.11 e cocotb 2.0.1.
  \item A AURORA não usa o cocotb pela linha de comando dele. Ela gera um script
        runner e passa os dados do projeto por variáveis de ambiente.
  \item Segurança: só o Python embarcado pode ser executado pelo fluxo de
        compilação; qualquer outro é rejeitado (Seção 7).
\end{itemize}

\section{O bundle MSYS2}

O Python não é instalado à parte. Ele vem dentro de um único bundle que contém
toda a toolchain FOSS empacotada a partir de um snapshot do MSYS2/mingw64.

\begin{itemize}
  \item Arquivo: \texttt{aurora-msys-v1.zip}, tag de release \texttt{msys-v1}.
  \item Origem: repositório \texttt{nipscernlab/aurora-toolchain} no GitHub
        (\url{https://github.com/nipscernlab/aurora-toolchain}). O download é
        feito pelo script \texttt{components/Scripts/download-toolchain.js}
        durante o \texttt{npm run bootstrap}.
  \item Conteúdo relevante: \texttt{mingw64/bin} (iverilog, vvp, yosys, verilator,
        perl, g++, make, ccache, e o \texttt{python.exe} com suas DLLs),
        \texttt{mingw64/lib/python3.12} (a stdlib e o \texttt{site-packages} com o
        cocotb), e \texttt{usr/bin} (bash e coreutils).
\end{itemize}

O cocotb não é instalado com \texttt{pip} em tempo de execução. Ele já vem
pré-montado dentro do zip, pela CI do \texttt{aurora-toolchain}, que fixa o MSYS2,
monta o VPI do cocotb, monta o bundle, faz um smoke test dos fluxos e publica.

A instalação do bundle é verificada por sentinelas: o \texttt{verilator\_bin.exe}
confirma a toolchain, e a presença de
\texttt{site-packages/cocotb/libs/libcocotbvpi\_verilator.a} confirma que o cocotb
veio junto. Um bundle sem cocotb dispara um novo download.

Links úteis:

\begin{itemize}
  \item cocotb: \url{https://www.cocotb.org} e \url{https://docs.cocotb.org}
  \item MSYS2: \url{https://www.msys2.org}
  \item aurora-toolchain (nosso empacotamento):
        \url{https://github.com/nipscernlab/aurora-toolchain}
\end{itemize}

\section{O interpretador e as versões}

Quem resolve qual Python usar é \texttt{main/compile/python\_locator.js}. O
interpretador canônico é o embarcado:

\begin{lstlisting}[language=bash]
components/Packages/msys/mingw64/bin/python.exe
\end{lstlisting}

O \texttt{python\_locator.js} monta uma lista de candidatos, na ordem: o Python
embarcado, depois as variáveis \texttt{AURORA\_PYTHON}, \texttt{PYTHON},
\texttt{CONDA\_PREFIX}, instalações comuns do Windows (miniconda, anaconda,
Python da python.org, \texttt{py.exe}) e por fim o que estiver no \texttt{PATH}.
Cada candidato é validado por uma sonda que roda um \texttt{python -c} curto e
reporta a versão do Python, se o cocotb está presente e qual a versão do cocotb.
Na prática o embarcado sempre vence: se ele existe, é o escolhido.

Versões fixadas:

\begin{itemize}
  \item Python 3.12.11 (a stdlib fica em \texttt{mingw64/lib/python3.12}).
  \item cocotb 2.0.1, fixado pela constante \texttt{EXPECTED\_COCOTB\_VERSION} em
        \texttt{python\_locator.js}. O status do Python (exposto pela IPC
        \texttt{toolchain:python-status}) marca se a versão do cocotb bate com a
        esperada.
\end{itemize}

\section{Bibliotecas instaladas}

Este é o ponto mais importante para entender o estado atual. O nosso Python é
mínimo. O \texttt{site-packages} contém apenas o cocotb e o que ele precisa:

\begin{itemize}
  \item \texttt{cocotb} (2.0.1): o framework.
  \item \texttt{cocotb\_tools}: o runner e a configuração do cocotb (o
        \texttt{runner.py} que a AURORA usa).
  \item \texttt{pygpi}: a ponte do cocotb com o simulador (a interface GPI).
  \item \texttt{find\_libpython} (0.5.1): dependência do cocotb.
\end{itemize}

E, igualmente importante, o que NÃO está instalado:

\begin{itemize}
  \item Não há \texttt{pip} nem \texttt{ensurepip}. O bundle não traz gerenciador
        de pacotes, então não dá para \texttt{pip install} algo direto de fábrica.
  \item Não há \texttt{setuptools} nem \texttt{wheel}.
  \item Não há bibliotecas científicas (\texttt{numpy}, \texttt{scipy}, etc.). Um
        testbench que faça análise numérica precisa delas e hoje elas não existem
        no runtime embarcado.
\end{itemize}

Ou seja, hoje o runtime cobre exatamente o fluxo cocotb e nada além. É esse vazio
que o repositório de bibliotecas (Seção 8) pretende preencher.

\section{Como um testbench cocotb roda}

O fluxo, do clique ao resultado, é o seguinte.

\paragraph{Detecção.} A AURORA decide pelo nome do arquivo. Um testbench
\texttt{.py} (função \texttt{isPythonFile} em
\texttt{js/compilation/compilation\_helpers.ts}) é roteado para o fluxo cocotb;
um \texttt{.v/.sv} vai para o fluxo iverilog/verilator tradicional. Isso vale
tanto para o botão Wave (\texttt{runGtkWave}) quanto para o Fast Sim
(\texttt{runFastSim}), ambos em
\texttt{js/compilation/compilation\_module.js}.

\paragraph{Módulo topo (DUT).} Dentro do \texttt{.py}, um comentário opcional
indica o módulo HDL de topo:

\begin{lstlisting}[language=python]
# aurora-toplevel: meu_modulo
\end{lstlisting}

Se presente, esse é o DUT. Se ausente, cai no top-level do projeto \texttt{.spf}.
Por ser um comentário, o cocotb puro ignora a linha.

\paragraph{O runner gerado.} A AURORA não chama o cocotb pela CLI. Ela gera, em
tempo de execução, o script \texttt{components/Temp/aurora\_cocotb\_runner.py}
(método \texttt{\_writeCocotbRunnerScript}) e o executa. Esse runner lê as
variáveis de ambiente, injeta os caminhos no \texttt{sys.path} e usa a API
oficial do cocotb (\texttt{get\_runner} de \texttt{cocotb\_tools.runner}) para
compilar o HDL (\texttt{runner.build}) e rodar o testbench (\texttt{runner.test}).
O comando executado é, em essência:

\begin{lstlisting}[language=bash]
python.exe aurora_cocotb_runner.py
\end{lstlisting}

com tudo o mais vindo do ambiente.

\paragraph{O contrato de ambiente.} Os dados do projeto vão por variáveis
\texttt{AURORA\_COCOTB\_*}. As quatro protegidas (que o override da Aurora
Intelligence não pode alterar, ver \texttt{main/compile/protected\_flags.js}):

\begin{lstlisting}
AURORA_COCOTB_SOURCES_JSON   # JSON com os arquivos HDL passados ao runner.build
AURORA_COCOTB_TOP            # nome do modulo topo (hdl_toplevel, o DUT)
AURORA_COCOTB_TEST_MODULE    # nome do modulo Python (o .py) a importar
AURORA_COCOTB_BUILD_DIR      # diretorio de build/saida (Temp/cocotb_<tb>)
\end{lstlisting}

Além dessas, o ambiente carrega o diretório de teste, o \texttt{PYTHONPATH} (para
o \texttt{.py} e os módulos do projeto importarem), os argumentos de build por
engine, o \texttt{SIM} (\texttt{icarus} ou \texttt{verilator}), o
\texttt{TOPLEVEL\_LANG}, o \texttt{WAVES} (1 no Wave, 0 no Fast Sim) e ajustes de
encoding (\texttt{PYTHONUTF8}, \texttt{PYTHONIOENCODING}) para o log não quebrar.

\paragraph{A execução.} O runner é lançado pelo processo principal
(\texttt{main/compile/executor.js}) com \texttt{shell:false} (sem shell, sem
injeção), registrado no process registry para que o cancelamento derrube o Python
e tudo que ele ramifica (iverilog/verilator/g++). O stdout aparece ao vivo no
terminal.

\paragraph{Engines.} O mesmo Python roda os dois engines; muda só o \texttt{SIM}
e os argumentos de build. Icarus usa \texttt{-g2012}. Verilator usa um conjunto
de \texttt{-Wno-*} (para os warnings do HDL do SAPHO não abortarem o build), o
\texttt{+define+YANC\_TRACE} e flags de otimização. No modo Wave o Verilator
recebe \texttt{--trace-fst} e \texttt{WAVES=1}; no Fast Sim roda headless, sem
onda.

\paragraph{Progresso e resultado.} A AURORA lê o stdout do cocotb procurando o
padrão de progresso (regex \texttt{COCOTB\_PROG\_RE} em
\texttt{compilation\_module.js}, algo como \texttt{"nome: N/M samples
processed"}) e transforma isso na barra de progresso do terminal. O sucesso ou
falha é dado pelo código de saída do processo: o \texttt{runner.test} retorna
diferente de zero quando os testes falham. No modo Wave, o dump \texttt{.fst}/
\texttt{.vcd} gerado é localizado e aberto no viewer.

\section{Onde fica cada componente}

Agrupando pelos papéis, com os nomes reais:

\paragraph{Localizar e validar o interpretador.}
\texttt{main/compile/python\_locator.js} (resolve o Python embarcado, descobre
candidatos, sonda o cocotb, expõe \texttt{getPythonStatus}),
\texttt{main/paths.js} (a raiz \texttt{componentsPath} do bundle),
\texttt{main/ipc/system.js} (a IPC \texttt{toolchain:python-status}).

\paragraph{Bundle e download.}
\texttt{components/Scripts/download-toolchain.js} (baixa e extrai o
\texttt{aurora-msys-v1.zip}, checa a sentinela do cocotb). O interpretador e o
cocotb em si ficam em \texttt{components/Packages/msys/mingw64/}.

\paragraph{Montar e rodar o comando cocotb.}
\texttt{js/compilation/builders/cocotb.ts} (monta o CommandSpec
\texttt{cocotb-run}), \texttt{js/compilation/compilation\_module.js} (a
orquestração: gera o runner, coleta as fontes, resolve o perfil Icarus/Verilator,
roda a simulação, adota a onda), \texttt{js/compilation/compilation\_helpers.ts}
(\texttt{isPythonFile}, a diretiva \texttt{aurora-toplevel}, validação de nome de
módulo), \texttt{js/wave/simulator\_preference.js} (o toggle do engine).

\paragraph{Executar (processo principal).}
\texttt{main/compile/executor.js} (spawn com \texttt{shell:false}, ambiente
sanitizado, prioridade elevada, slot de cancelamento) e
\texttt{main/process\_registry.js} (o tree-kill do grupo RUN).

\paragraph{Segurança.}
\texttt{main/compile/binary\_allowlist.js} (só o Python embarcado é executável) e
\texttt{main/compile/protected\_flags.js} (as invariantes de ambiente do
\texttt{cocotb-run}).

\section{Segurança}

Dois controles cercam a execução do Python:

\begin{itemize}
  \item Allowlist de binário: em \texttt{binary\_allowlist.js}, qualquer nome
        \texttt{python}/\texttt{python3}/\texttt{py} só é liberado se o caminho
        for exatamente o do interpretador embarcado. Um Python de fora é
        recusado, não importam os argumentos. É o mesmo mecanismo que restringe
        os outros binários da toolchain.
  \item Flags protegidas: o spec \texttt{cocotb-run} não carrega flags de CLI
        sensíveis (só o caminho do runner), então as invariantes são as quatro
        variáveis \texttt{AURORA\_COCOTB\_*} da Seção 5. Um override não pode
        mudar o valor delas.
  \item \texttt{shell:false} no spawn: os argumentos vão diretos ao processo, sem
        passar por um shell, o que elimina injeção por parsing de string.
\end{itemize}

\section{O painel de bibliotecas}

O vazio da Seção 4 é preenchido por um painel dentro da AURORA que lista as
bibliotecas disponíveis, com descrição e casos de uso, e instala sob demanda. As
bibliotecas que não couberem nesse caminho ficam a cargo do Python do próprio
usuário, via terminal TCMD.

\subsection{A restrição técnica que orienta tudo}

O nosso Python é um build MinGW do MSYS2, não o CPython oficial (MSVC) da
python.org. Três testes feitos diretamente sobre o interpretador embarcado
delimitam o que é possível:

\begin{itemize}
  \item Biblioteca Python pura (wheel \texttt{*-none-any.whl}): funciona. O
        Plotly 6.9.0 da PyPI importou e gerou figura HTML sem nenhum ajuste.
  \item Wheel binária da PyPI (\texttt{cp312-win\_amd64}): não carrega. Mesmo
        forçando o nome do \texttt{.pyd}, o erro é \texttt{DLL load failed} —
        o módulo procura a \texttt{python312.dll} do CPython da Microsoft, que
        não existe no nosso build, que expõe \texttt{libpython3.12.dll}.
  \item Pacote mingw do MSYS2: também não. Este ponto corrige o plano anterior,
        que apostava em \texttt{mingw-w64-x86\_64-python-numpy} por ser compilado
        para o mesmo compilador. O raciocínio de ABI está certo, mas o MSYS2 é
        \emph{rolling release} e já andou: o repositório \texttt{mingw64} hoje
        entrega Python 3.14.6, e o NumPy 2.4.6-1 de lá vem marcado \texttt{cp314},
        instalado em \texttt{lib/python3.14/}. Contra o nosso 3.12.11 o import
        falha. Ou seja, a barreira imediata não é o compilador: é a versão do
        Python. O bundle está congelado num snapshot de junho de 2025 e o
        repositório do MSYS2 só distribui o presente.
\end{itemize}

O critério prático que decorre disso é verificável por máquina, e não por
julgamento: se a PyPI publica uma wheel terminada em \texttt{-none-any.whl}, a
biblioteca instala; se só publica wheel compilada, não instala de jeito nenhum.
O painel consulta isso na API da PyPI \emph{antes} de baixar qualquer byte, tanto
para a lista curada quanto para uma biblioteca arbitrária pedida pelo usuário.

\subsection{Arquitetura}

O catálogo (\texttt{resources/pylib-catalog.json}) guarda metadados, não bytes:
nome, versão fixada, descrição e casos de uso em português e inglês, licença,
categoria, ícone e o \texttt{sha256} de cada wheel. Ele é gerado por
\texttt{scripts/gen-pylib-catalog.js}, que consulta a PyPI e recusa sozinho
qualquer entrada que não seja pura. As bibliotecas compiladas entram como
registros informativos, sem wheel, para o painel poder exibi-las e explicar o
motivo em vez de fingir que não existem.

Os bytes vêm da PyPI no momento da instalação e são conferidos contra o hash
fixado. Não há resolvedor de dependências, e isso é proposital: sem \texttt{pip}
no runtime não existe resolvedor, então o catálogo já traz o fecho completo
pronto e o instalador apenas obedece à lista.

\subsection{Para que servem, e onde ficam}

As bibliotecas instaladas pelo painel valem para \emph{qualquer} arquivo Python
executado pelo interpretador embarcado: um testbench cocotb, um script solto do
projeto, uma linha digitada no terminal. Não são um recurso do fluxo cocotb.

\subsection{Onde as bibliotecas são instaladas}

Em \texttt{components/PyLibs/site/}, e não no \texttt{site-packages} do bundle.
A pasta do bundle pertence à toolchain: o bootstrap a baixa, o
\texttt{scripts/verify-components.js} a re-baixa quando o componente está
faltando ou desatualizado, e o \texttt{components/Packages/} inteiro é
\emph{gitignored} por ser artefato derivado. Instalar estado do usuário lá dentro
significaria perdê-lo num \texttt{--force} ou quando alguém apaga a pasta para
consertar a toolchain, que é o conserto documentado.

\subsection{Como o interpretador encontra as bibliotecas}

Sem \texttt{pip} nem \texttt{ensurepip}, \texttt{venv} não é uma opção real. A
ligação é feita por um arquivo \texttt{aurora-pylibs.pth} escrito dentro do
\texttt{site-packages} do próprio interpretador embarcado. O Python lê, na
inicialização, todo \texttt{.pth} que encontra ali e acrescenta ao
\texttt{sys.path} as pastas listadas. Isso resolve os dois requisitos de uma vez:

\begin{itemize}
  \item \textbf{Alcance.} Vale para qualquer execução deste interpretador, sem
        depender de quem o chamou nem de exportar variável de ambiente.
  \item \textbf{Isolamento.} Vale só para este interpretador. O \texttt{site.py}
        de cada Python lê apenas os \texttt{.pth} do \texttt{site-packages} dele,
        então o Python que o usuário tem instalado na máquina nunca enxerga o
        nosso \texttt{PyLibs} e não há colisão com o que ele instalou por
        \texttt{pip}.
\end{itemize}

Uma variável \texttt{PYTHONPATH} no ambiente do terminal faria o oposto: seria
lida por qualquer Python que rodasse ali, inclusive o do usuário. Houve uma versão
que acrescentava o diretório ao \texttt{AURORA\_COCOTB\_PYTHONPATH} em
\texttt{main/compile/executor.js}; funcionava, mas só para o cocotb, e foi
substituída pelo \texttt{.pth} para não manter dois mecanismos com o mesmo fim. O
\texttt{AURORA\_COCOTB\_PYTHONPATH} segue nas chaves protegidas do passo
\texttt{cocotb-run}, porque continua sendo um caminho de carregamento de código.

O arquivo é reescrito na abertura do app (\texttt{main/python/pylib\_watch.js}),
porque uma reinstalação da toolchain o apaga junto com o bundle. Ele não é estado:
é um ponteiro de uma linha, e nada se perde se desaparecer.

\subsection{Escolher o Python no terminal}

O TCMD é um shell de verdade, e \texttt{python} ali significa o Python que o
usuário instalou. Isso não lhe é tomado. A sessão ganha três comandos, definidos
no mesmo \emph{bootstrap} que carrega o \emph{prompt}, sem tocar em
\texttt{PATH} nem em \texttt{PYTHONPATH}:

\begin{itemize}
  \item \texttt{apython} — invoca sempre o interpretador embarcado.
  \item \texttt{Use-Python aurora} — faz \texttt{python} significar o nosso nesta
        sessão, definindo uma função com esse nome (em PowerShell, função tem
        precedência sobre executável do \texttt{PATH}).
  \item \texttt{Use-Python system} — apaga a função, e \texttt{python} volta a ser
        o do usuário.
\end{itemize}

O padrão é \texttt{system}: o terminal começa sendo o shell que o usuário espera,
e a troca é um ato explícito dele.

\subsection{Instalar, remover, reparar}

Instalar é baixar a wheel medindo o progresso real em bytes, conferir o
\texttt{sha256}, recusar qualquer caminho interno com \texttt{..} ou absoluto,
extrair e anotar no manifesto (\texttt{components/PyLibs/installed.json}) a lista
completa de arquivos escritos. É essa anotação que torna a remoção exata.

Remover apaga apenas os arquivos que nenhuma outra biblioteca instalada
reivindica — o \texttt{packaging} chega junto do Plotly e do pytest, o
\texttt{pygments} junto do pytest e do Rich — e em seguida varre os diretórios de
topo que ficaram sem dono, levando junto o \texttt{\_\_pycache\_\_} que o Python
criou depois da instalação e que não existia na wheel.

Reparar é remover e instalar de novo. Atende ao caso em que o manifesto diz
``instalada'' mas os arquivos sumiram, e ao caso de acompanhar uma versão nova do
catálogo. O manifesto grava contra qual ABI a instalação foi feita, de modo que
um futuro bump do Python do bundle seja detectado e avisado, em vez de quebrar em
silêncio — que é exatamente o que aconteceu com o pacote mingw64 do NumPy.

\subsection{O que ainda falta: as bibliotecas compiladas}

NumPy, SciPy, Matplotlib e pandas não instalam por esse caminho, e não há ajuste
no instalador que resolva: é incompatibilidade de ABI, não de configuração.

Isso pesa mais do que parece, porque é justamente o ferramental de que a
metodologia publicada do grupo depende. O testbench cocotb da PMU segundo a norma
IEC/IEEE 60255-118-1 usa geração de sinal, modelo de referência em dupla
precisão, filtro FIR classe M e estatística de TVE, FE e ROCOF; o trabalho da CNN
do TileCal compara a saída do SAPHO com a referência em Python por erro médio
absoluto e $R^2$.

O que o painel entrega hoje cobre modelagem de referência e conferência, mas não
substitui o NumPy em volume: \texttt{mpmath} permite fixar a mantissa em 24 bits
e modelar o \emph{float} próprio do SAPHO (1+8+23, fora do IEEE 754), separando
erro de método de erro de arredondamento; \texttt{sympy} deduz coeficientes de
filtro de forma exata; \texttt{vcdvcd} e \texttt{pyvcd} leem e escrevem forma de
onda; o Plotly desenha a conferência. Um laço de filtragem sobre milhares de
amostras em Python puro, porém, é lento.

Há três saídas, em ordem de custo:

\begin{itemize}
  \item Levar as compiladas para dentro do \textbf{bundle}. O
        \texttt{aurora-toolchain} já monta o bundle a partir de um snapshot
        fixado do MSYS2; NumPy e SciPy podem entrar pela mesma CI, do mesmo
        snapshot, com a ABI garantida por construção. Custo: o bundle cresce (o
        NumPy sozinho ocupa cerca de 55\,MB descompactado e arrasta o
        \texttt{openblas}) e cada mudança vira um bump de toolchain.
  \item Subir o Python do bundle para 3.14 e acompanhar o MSYS2, passando a usar
        os pacotes mingw prontos. Custo: mexe no cocotb, no VPI do Verilator e na
        toolchain inteira de uma vez.
  \item Deixar com o usuário, via TCMD e o Python dele. Custo zero para nós, mas
        o modelo de referência sai de dentro do fluxo da AURORA.
\end{itemize}

A primeira é a recomendada: o painel continua simples e sem risco de ABI, as
bibliotecas puras chegam sob demanda, e as pesadas viram decisão de bundle, rara
e verificada por CI.

\subsection{Fallback: o Python do usuário via TCMD}

O TCMD é o terminal interativo da AURORA, um pseudo-terminal real (ConPTY, via
\texttt{@lydell/node-pty}) que roda um shell de verdade (PowerShell no Windows).
Nele o usuário roda comandos livres, incluindo o Python que ele tiver instalado
na máquina:

\begin{lstlisting}[language=bash]
python meu_script.py
\end{lstlisting}

Esse é o escape para as bibliotecas que não estiverem no nosso repositório: em
vez de forçar tudo para dentro do runtime embarcado (com a restrição de ABI), o
usuário usa o Python dele, com o ecossistema completo da PyPI, direto no TCMD. O
runtime embarcado continua dedicado ao cocotb; o TCMD cobre o uso geral.

\end{document}
